\section{Sample results and economic reflection}
\label{sec:results}

This section reports a representative backtest, contextualises it against the two
benchmark overlays the system computes automatically, and reflects on what the
result implies about the strategy's regime sensitivity. The numeric cells in
Table~\ref{tab:results} are populated by the team from a run they execute
themselves, so the table represents the actual code path rather than a fabricated
example.

\subsection{Run methodology}

\paragraph{Configuration.} We report a backtest of the \texttt{sma-crossover}
strategy on \texttt{SPY} (SPDR S\&P~500 ETF Trust), over 2020-01-01 to 2024-12-31,
daily bars (\texttt{timeframe="1d"}). Strategy parameters:
\texttt{fast\_window=20}, \texttt{slow\_window=50}, \texttt{allow\_short=false}.
Execution parameters: \texttt{initial\_cash=10\,000}, \texttt{commission\_bps=5},
\texttt{slippage\_bps=2}, \texttt{sizing\_mode="fixed\_fraction"},
\texttt{risk\_fraction=1.0}, no max-DD circuit breaker. The annualisation factor
is $k=252$ (NYSE daily, Section~\ref{sec:calendars}).

\paragraph{Reproducibility.} The exact API call:

\begin{lstlisting}[language=bash]
curl -s -X POST http://127.0.0.1:8000/backtests \
  -H 'Content-Type: application/json' \
  -d '{
    "asset_symbol": "SPY",
    "strategy_slug": "sma-crossover",
    "start_date":   "2020-01-01T00:00:00",
    "end_date":     "2024-12-31T00:00:00",
    "params":       {"fast_window": 20, "slow_window": 50, "allow_short": false},
    "timeframe":    "1d"
  }'
\end{lstlisting}

Then \texttt{GET /backtests/\{run\_id\}/metrics} and the two
\texttt{/benchmark/\{kind\}/metrics} endpoints supply the table values.

\subsection{Results table}

\todo[inline]{TEAM: replace the \emph{TBD} cells in Table~\ref{tab:results} below
with the numbers your actual run returned. The table compiles either way; the
TBDs make the gap visible to a reviewer.}

\begin{table}[ht]
\centering
\small
\renewcommand{\arraystretch}{1.2}
\begin{tabular}{lccc}
\toprule
\textbf{Metric} & \textbf{SMA Crossover} & \textbf{Same-asset B\&H} & \textbf{SPY B\&H} \\
\midrule
Total return       & \textit{TBD} & \textit{TBD} & \textit{TBD} \\
CAGR               & \textit{TBD} & \textit{TBD} & \textit{TBD} \\
Annualised vol     & \textit{TBD} & \textit{TBD} & \textit{TBD} \\
Sharpe             & \textit{TBD} & \textit{TBD} & \textit{TBD} \\
Sortino            & \textit{TBD} & \textit{TBD} & \textit{TBD} \\
Max drawdown       & \textit{TBD} & \textit{TBD} & \textit{TBD} \\
Calmar             & \textit{TBD} & \textit{TBD} & \textit{TBD} \\
Trade count        & \textit{TBD} & 1 & 1 \\
Win rate           & \textit{TBD} & --- & --- \\
Profit factor      & \textit{TBD} & --- & --- \\
\bottomrule
\end{tabular}
\caption{KPIs for the representative SMA-crossover run, alongside the two
benchmark overlays computed automatically by \texttt{BacktestAgent}. Trade-block
cells are dashed for benchmarks (buy-and-hold has only one entry and one exit).
Cells marked \emph{TBD} are filled in by the team from the actual run.}
\label{tab:results}
\end{table}

\subsection{Economic reflection}

The economic story we expect Table~\ref{tab:results} to tell follows from the
nature of the strategy and the chosen window. SMA Crossover is a trend-follower:
it goes long when the 20-day SMA crosses above the 50-day SMA and exits when it
crosses back. The 2020\textendash{}2024 window spans the COVID drawdown (March
2020), the equity bull market of 2020\textendash{}2021, the 2022 bear market, and
the 2023\textendash{}2024 recovery \textemdash{} four distinct regimes inside five
years. A 20/50 SMA crossover should capture the trends of 2021 and 2023\textendash{}2024
cleanly, sit in cash during much of 2022, but also lose money to whipsaws around
the regime transitions (early 2020 and mid-2022).

We therefore expect three observations to hold once the team populates the table:

\begin{enumerate}
  \item \textbf{Lower CAGR than SPY buy-and-hold} but \textbf{noticeably lower
        maximum drawdown}, because the strategy goes flat during 2022's bear
        market. The pain-adjusted comparison (Calmar) is the more interesting
        head-to-head.
  \item \textbf{A Sharpe ratio close to buy-and-hold's} \textemdash{} the
        strategy's higher hit rate during clear trends is offset by the
        cost of being in cash through the recoveries.
  \item \textbf{A trade count modest enough to demand caution.} Five years of daily
        bars typically produce 10\textendash{}30 trades for a 20/50 SMA system, which
        is below the rule-of-thumb threshold of $\sim$100 trades for statistical
        significance. Any strong single-asset result here should be re-tested on
        several assets before being trusted.
\end{enumerate}

\todo[inline]{TEAM: replace the three-point list above with the actual observations
once Table~\ref{tab:results} is populated. The numbered list is a template; rewrite
it in your own words to reflect the run you actually executed.}

\subsection{Limitations of the demonstration}

The result above is a single asset over a single window with a single parameter
set. It is illustrative, not predictive. Three caveats apply:

\begin{itemize}
  \item \textbf{Survivorship bias} \textemdash{} yfinance carries only currently
        listed tickers, so any equity universe is biased upward.
  \item \textbf{Single-window evaluation} \textemdash{} the rubric's preference for
        walk-forward / out-of-sample evaluation is correct; SMA-crossover's
        in-sample fit on a single window does not establish edge.
  \item \textbf{The README's retail-trader tip applies} \textemdash{} if the
        backtested maximum drawdown is 20\%, the user should mentally prepare for
        30\% in live trading, both because of regime change and because realistic
        slippage on retail accounts is usually worse than the 2-bps default.
\end{itemize}
